%; whizzy paragraph -pdf xpdf -latex ./whizzypdfptex.sh
%; whizzy-paragraph "^\\\\begin{frame}\\|\\\\emtext"
% latex beamer presentation.
% platex, latex-beamer でコンパイルすることを想定。 

%     Tokyo Debian Meeting resources
%     Copyright (C) 2012 Junichi Uekawa

%     This program is free software; you can redistribute it and/or modify
%     it under the terms of the GNU General Public License as published by
%     the Free Software Foundation; either version 2 of the License, or
%     (at your option) any later version.

%     This program is distributed in the hope that it will be useful,
%     but WITHOUT ANY WARRANTY; without even the implied warreanty of
%     MERCHANTABILITY or FITNESS FOR A PARTICULAR PURPOSE.  See the
%     GNU General Public License for more details.
%     You should have received a copy of the GNU General Public License
%     along with this program; if not, write to the Free Software
%     Foundation, Inc., 51 Franklin St, Fifth Floor, Boston, MA  02110-1301 USA

\documentclass[cjk,dvipdfmx,12pt]{beamer}
\usetheme{Tokyo}
\usepackage{monthlypresentation}

%  preview (shell-command (concat "evince " (replace-regexp-in-string "tex$" "pdf"(buffer-file-name)) "&")) 
%  presentation (shell-command (concat "xpdf -fullscreen " (replace-regexp-in-string "tex$" "pdf"(buffer-file-name)) "&"))
%  presentation (shell-command (concat "evince " (replace-regexp-in-string "tex$" "pdf"(buffer-file-name)) "&"))

%http://www.naney.org/diki/dk/hyperref.html
%日本語EUC系環境の時
\AtBeginDvi{\special{pdf:tounicode EUC-UCS2}}
%シフトJIS系環境の時
%\AtBeginDvi{\special{pdf:tounicode 90ms-RKSJ-UCS2}}

\newenvironment{commandlinesmall}%
{\VerbatimEnvironment
  \begin{Sbox}\begin{minipage}{1.0\hsize}\begin{fontsize}{8}{8} \begin{BVerbatim}}%
{\end{BVerbatim}\end{fontsize}\end{minipage}\end{Sbox}
  \setlength{\fboxsep}{8pt}
% start on a new paragraph

\vspace{6pt}% skip before
\fcolorbox{dancerdarkblue}{dancerlightblue}{\TheSbox}

\vspace{6pt}% skip after
}
%end of commandlinesmall

\title{東京エリアDebian勉強会}
\subtitle{第155回 2017年9月度}
\author{杉本 典充  dictoss@live.jp}
\date{2017年9月16日}
\logo{\includegraphics[width=8cm]{image200607/openlogo-light.eps}}

\begin{document}

\begin{frame}
\titlepage{}
\end{frame}

\begin{frame}{Agenda}
 \begin{minipage}[t]{0.45\hsize}
  \begin{itemize}
  \item 最近あったDebian関連のイベント報告
    \begin{itemize}
    \item 第153回東京エリアDebian勉強会
    \item 第154回東京エリアDebian勉強会 (OSC 2017 Tokyo/Fall)
    \end{itemize}
  \item 事前課題発表
  % \item Debian Trivia Quiz
  \end{itemize}
 \end{minipage}
 \begin{minipage}[t]{0.45\hsize}
  \begin{itemize}
  \item セミナー発表
    \begin{itemize}
    \item 初めてのキーサインパーティ
    \item Debian on Pomera DM200
    \end{itemize}
  \end{itemize}
 \end{minipage}
\end{frame}

\section{イベント報告}
\emtext{イベント報告}

\begin{frame}{第153回東京エリアDebian勉強会}
\begin{itemize}
\item 2017年8月19日に、朝日ネットさんをお借りして開催
\item 参加者は9名
\item セミナー「debconf17参加報告」(青木さん、やまねさん）
\end{itemize} 
\end{frame}


\begin{frame}{第154回東京エリアDebian勉強会 (OSC 2017 Tokyo/Fall)}
\begin{itemize}
\item 2017年9月10日（2日目のみ）に、明星大学様で開催
\item 2017年9月10日のイベント来場者数は、約450名
\item セミナー参加者は、10数名
\item セミナー「Debian Updates」(dictossさん）
\item 展示ブースへの来訪は、約45名
\item Debianをインストールする小学生
\item Raspbian使いの方は多かったが、Debianは使わずRaspberry Pi実機のみを使うという方がそこそこ存在することに印象を持った
\end{itemize} 
\end{frame}

\section{事前課題}
\emtext{事前課題}
{\footnotesize
 %
% $BDs=P2]Bj(B
%
\begin{prework}{ ichinomoto }
  \begin{enumerate}
  \item DM200$B$N(BTODO$B$K$J$C$F$$$kItJ,$ND4::$J$I(B
  \item systemd$B<~$j$N@_Dj$K$D$$$F(B
  \end{enumerate}
\end{prework}

\begin{prework}{ dictoss }
  \begin{enumerate}
  \item KDE$B$H(Blibqt5$B$N;H$$J}$NJY6/(B
  \item debian.org$B$H(Bdebian.net$B$G?J$s$G$$$k%W%m%8%'%/%H0lMw$H$=$N>u67(B
  \end{enumerate}
\end{prework}

\begin{prework}{ ysaito }
  \begin{enumerate}
  \item $B<!(B,$B2?$+H/I=$G$-$k$b$N$,$"$k$+9M$($^$9!#(B
  \item ($B2sEz$J$7(B)
  \end{enumerate}
\end{prework}

\begin{prework}{ henrich }
  \begin{enumerate}
  \item $B;(CL$+$b!D!#(B
  \item $B;H$C$F$$$F6qBNE*$K:$$C$F$$$k$3$H!#(B
  \end{enumerate}
\end{prework}

\begin{prework}{ koedoyoshida }
  \begin{enumerate}
  \item $BL$Dj(B
  \item ($B2sEz$J$7(B)
  \end{enumerate}
\end{prework}

\begin{prework}{ NOKUBI Takatsugu }
  \begin{enumerate}
  \item uim/anthy$B$NLdBj2r7h!"(Bmlterm$B$N(Bsixel$BM-8z2=(B
  \item $B5!3#3X=,%M%?$rH/I=$G$-$J$$$+$J$H9M$($F$$$^$9(B
  \end{enumerate}
\end{prework}

\begin{prework}{ Kouhei Maeda }
  \begin{enumerate}
  \item $B%Q%C%1!<%8%a%s%F%J%s%9(B
  \item ($B2sEz$J$7(B)
  \end{enumerate}
\end{prework}

}

%\subsection{問題}
%%; whizzy-master ../debianmeetingresume201311.tex
% $B0J>e$N@_Dj$r$7$F$$$k$?$a!"$3$N%U%!%$%k$G(B M-x whizzytex $B$9$k$H!"(Bwhizzytex$B$,MxMQ$G$-$^$9!#(B
%

\santaku
{quetion$B!#(B}
{A}
{B}
{C}
{C}
{desc}


\section{セミナー1}
\emtext{初めてのキーサインパーティ}

\section{セミナー2}
\emtext{Debian on Pomera DM200}

\section{Hack time}
\emtext{Hack time}

%\begin{frame}{成果を記入下さい!}
%\end{frame}
  
\section{今後のイベント}
\emtext{今後のイベント}
\begin{frame}{今後のイベント}
  \begin{itemize}
  \item 2017/10/21(土) 東京エリアDebian勉強会
    \begin{itemize}
    \item 13:00-18:00で、会場は朝日ネット様の予定
    \item セミナー内容は調整中
    \end{itemize}
  \item 2017/11/18(土) 東京エリアDebian勉強会 (予定)
    \begin{itemize}
    \item 次回、場所、セミナー内容は調整中
    \end{itemize}    
  \item 勉強会では発表者を随時募集中です。ぜひともよろしくお願いします。
  \end{itemize}
\end{frame}

\end{document}

;;; Local Variables: ***
;;; outline-regexp: "\\([ 	]*\\\\\\(documentstyle\\|documentclass\\|emtext\\|section\\|begin{frame}\\)\\*?[ 	]*[[{]\\|[]+\\)" ***
;;; End: ***
